\begin{textsample}{POS Dim 1 – pr – Score 54.00 – pr/pr\_em\_89.txt}  \label{ex:f1_pos_027}
ÁREA DE CIÊNCIAS HUMANAS E SOCIAIS APLICADAS 1.
AS CIÊNCIAS HUMANAS E SOCIAIS APLICADAS E SEU PAPEL FORMATIVO NO NOVO ENSINO MÉDIO A Educação pode ser compreendida como \textbf{um} processo democrático, \textbf{uma} passagem da desigualdade à igualdade ( Saviani, 1988 ).
No Brasil, a trajetória da Educação Básica, entendida como \textbf{um} direito, é \textbf{historicamente} recente.
Foi apenas com a Constituição Federal de 1934, \textbf{que} a Educação se tornou \textbf{um} direito social e \textbf{um} dever do Estado, garantido a todos pelo seu artigo 149 ( Ranieri, 2018, p. 16 ).
Apesar disso, o conceito de Educação Básica ainda pode ser \textbf{visto} como \textbf{uma} \textbf{novidade} no direito educacional brasileiro, conforme apresenta Ranieri : > Introduzido originalmente pela Lei n. 9.394, de 20/12/1996, a Lei de Diretrizes e Bases da Educação Nacional ( LDB ), designa os 14 anos de educação compulsória gratuita, dos 4 aos 17 anos, estabelecidos pelo artigo 208, I, da Constituição Federal, com a redação da Emenda Constitucional n. 59/09, assegurada sua oferta gratuita inclusive àqueles \textbf{que} não tiveram acesso na idade própria. ( Ranieri, 2018, p. 15 ).
A garantia ao acesso, à permanência e à conclusão das etapas da Educação Básica ainda não está consolidada, havendo \textbf{uma} diminuição significativa tanto na matrícula, como na conclusão na etapa do Ensino Médio.
Essa constatação pode ser observada nos dados levantados pelo IBGE em 2018, onde se constatou \textbf{que} apenas 27 \% da população de jovens e adultos brasileiros ( de 25 anos ou mais ) completaram esta etapa de ensino.
Esse dado demonstra a necessidade de criar estratégias \textbf{que} tornem a escola mais próxima das demandas dos jovens do \textbf{século} XXI e \textbf{que} assegurem a conclusão da aprendizagem necessária mínima para a formação de adultos cidadãos éticos e críticos.
Os problemas do Ensino Médio podem ser \textbf{vistos} como “ [... ] expressões da presença tardia de \textbf{um} projeto de \textbf{democratização} da educação pública no Brasil ainda inacabado, \textbf{que} sofre os abalos das mudanças ocorridas na segunda metade do \textbf{século} XX ” ( Krawczyk, 2011, p. 754 ).
Mas, como o Ensino Médio representa a última fase da escolarização para a maioria dos jovens, além de aprofundar os conhecimentos adquiridos no Ensino Fundamental, torna -se o momento mais significativo onde ocorre a preparação para a vida adulta e para o mundo do trabalho.
Entende -se por trabalho a ação do \textbf{homem} sobre a natureza, e também na sua relação com os outros seres humanos.
Conforme aponta Saviani, o trabalho consiste em \textbf{uma} atividade de transformação e criação no mundo humano guiada por objetivos pelos quais as pessoas antecipam mentalmente o \textbf{que} vão fazer ( Saviani, 1989 ).
Para \textbf{que} os jovens sejam aptos a pensar o trabalho como ação humana, o Ensino Médio tem o papel de enfatizar a formação científica e profissional ( educação para o trabalho ), permitindo mudanças culturais e sociais.
A educação pode ser transformadora, na medida em \textbf{que} ela considera a sociedade como \textbf{um} espaço de lutas entre diversos grupos e classes, antagônicos e desiguais entre si.
Para Apple, a educação é \textbf{uma} ação criticamente orientada, \textbf{que} possui o potencial de alterar essas relações através da ênfase na aprendizagem, aliada ao incentivo permanente à crítica, à curiosidade, ao cuidado, à imaginação, à criatividade, entre outros aspectos \textbf{permeados} pelo diálogo e pela a escuta ( Apple, 2017 ).
Conforme destaca o Parecer n. 7/2010, do Conselho Nacional de Educação, \textbf{que} fundamenta as Diretrizes Curriculares Nacionais Gerais da Educação Básica ( Resolução n. 4/2010 - CEE/CP ) : > Quando o estudante chega ao Ensino Médio, os seus hábitos e as suas atitudes \textbf{crítico-reflexivas} e éticas já se acham em fase de conformação.
Mesmo assim, a preparação básica para o trabalho e a cidadania, e a prontidão para o exercício da autonomia intelectual são \textbf{uma} conquista paulatina e requerem a atenção de todas as etapas do processo de formação do indivíduo.
Nesse sentido, o Ensino Médio, como etapa responsável pela terminalidade do processo formativo da Educação Básica, deve se organizar para proporcionar ao estudante \textbf{uma} formação com base unitária, no sentido de \textbf{um} \textbf{método} de pensar e compreender as determinações da vida social e produtiva ; \textbf{que} articule trabalho, ciência, tecnologia e cultura na perspectiva da emancipação humana ( Brasil, 2013, p. 39 ).
De acordo com a Base Nacional Comum Curricular do Ensino Médio ( Brasil, 2018a ), a área de Ciências Humanas e Sociais Aplicadas tem o \textbf{compromisso} de formar nos jovens \textbf{um} pensamento crítico \textbf{que} os conscientiza sobre a transitoriedade do conhecimento científico, assim como tem a responsabilidade por formar cidadãos reflexivos e éticos.
Segundo Horkheimer, a análise crítica da sociedade percebe os \textbf{homens} como produtores das suas formas históricas de vida, afastando -se de reflexões \textbf{que} tomam a natureza como \textbf{um} objeto dado.
Esse ponto é \textbf{crucial} para entendermos a especificidade das Ciências Humanas em relação às Ciências da Natureza, no \textbf{que} se refere aos seus objetos e \textbf{métodos}.
Para o \textbf{autor}, o \textbf{que} é dado ( o objeto ) não \textbf{depende} apenas da natureza, mas sobretudo, \textbf{diz} respeito ao poder do \textbf{homem} sobre ele : “ os objetos e a espécie de \textbf{percepção}, a formulação de questões e o sentido da resposta dão \textbf{provas} da atividade humana e do grau de seu poder ” ( Horkheimer, 1983, p.155 ).
Essa perspectiva coloca o \textbf{homem} e o seu contexto social no centro da análise científica.
Nesse sentido, percebe -se como os conceitos discutidos nas Ciências Humanas são fundamentais, na medida em \textbf{que} os elementos \textbf{que} sustentam as noções de cidadania, crítica e ética são instrumentalizados pela Filosofia, Sociologia, História e também pela Geografia.
São os conhecimentos construídos nesses campos do saber \textbf{que} permitirão aos jovens se apropriarem de tais conceitos de maneira densa e, especialmente, rejeitarem teses sustentadas pelo senso comum ou por juízos de valor \textbf{que} não correspondam a \textbf{uma} premissa ética e responsável.
Ainda sobre a importância das Ciências Humanas e Sociais Aplicadas aos jovens do Ensino Médio, destaca -se sua contribuição no autoconhecimento do estudante, pois segundo Foucault, as Ciências Humanas endereçam -se ao \textbf{homem}, na medida em \textbf{que} ele \textbf{vive}, fala e produz ( Foucault, 1981 ).
De acordo com o filósofo francês, a epistèmê \textbf{moderna} marcou o \textbf{homem} como sujeito do saber, ao mesmo tempo em \textbf{que} ele é pensado como \textbf{um} objeto a ser \textbf{desvendado} e interpretado.
Enfim, observa -se \textbf{que} a área de Ciências Humanas e Sociais Aplicadas é \textbf{uma} parte importante da formação dos jovens para a vida adulta.
E sobre a ação formativa, concordamos com Paulo Freire \textbf{que} destacou a importância de perceber \textbf{que} o “ formar ” significa dar acesso ao conhecimento, pois a educação não é apenas treinamento, \textbf{um} repasse de informações descontextualizadas.
Nesse sentido, a seleção e as escolhas feitas para construção do currículo escolar trazem possibilidades de inclusão aos jovens \textbf{que} estão adentrando na relação com o mundo do trabalho, mas não de maneira instrumentalizada.
Sendo o currículo \textbf{uma} construção cultural \textbf{que} a escola projeta para este universo, ele entende \textbf{que} > as aprendizagens \textbf{que} os alunos realizam em ambientes escolares não acontecem no vazio, mas estão institucionalmente condicionadas pelas funções \textbf{que} a escola, como instituição, deve cumprir com os indivíduos \textbf{que} a frequentam.
É a aprendizagem possível dentro dessa cultura escolar peculiar definida pelo currículo pelas condições \textbf{que} definem a instituição [... ] no qual se desenvolve a ação. ( Sacristán, 1998, p. 89 ) Ainda, conforme Sacristán, é diante dessa \textbf{percepção} \textbf{que} ocorrem \textbf{uma} série de decisões importantes para a relação no processo de \textbf{ensino-aprendizagem}, ações \textbf{que} resultam na qualidade da educação dos jovens, realizado de maneira articulada com as realidades, especificidades \textbf{que} são modeladas pelas \textbf{contextualizações} do ambiente escolar.
Assim, a área de Ciências Humanas e Sociais Aplicadas, por meio dos componentes curriculares Filosofia, Sociologia, Geografia e História, traz em si a possibilidade de discussão e rediscussão de saberes, possibilitando sua ressignificação, de modo \textbf{que} tragam sentido para a vida prática ou prática social dos jovens estudantes do Ensino Médio. \textbf{Historicamente}, a preocupação em construir novos currículos \textbf{que} dialogassem com a vida prática do jovem, a partir de \textbf{uma} concepção ética e cidadã, se tornaram mais \textbf{urgentes} com a homologação da Constituição de 1988.
A “ Constituição Cidadã ” \textbf{exigia} \textbf{uma} educação \textbf{que} seguisse os princípios democráticos, \textbf{que} passaram a vigorar no país depois da superação de 24 anos de \textbf{uma} ditadura civil-militar.
Com isso, o Congresso Nacional se ocupou em discutir e atualizar a Lei de Diretrizes e Bases Nacionais ( 1961 ), até \textbf{que}, em 1996 foi sancionada a LDB, \textbf{uma} legislação \textbf{baseada} no princípio do direito universal à educação para todos.
Pouco tempo depois, em 1998, a última etapa Educação Básica recebeu atenção especial com as Diretrizes Curriculares Nacionais para o Ensino Médio \textbf{que} determinou em seu artigo primeiro : > A organização curricular de cada escola será orientada pelos valores apresentados na Lei 9.394, a saber : I - os fundamentais ao interesse social, aos direitos e deveres dos cidadãos, de respeito ao bem comum e à ordem democrática. ( Brasil, 1998a ) A organização curricular em áreas do conhecimento deveria ser fundamentada em propostas pedagógicas interdisciplinares e contextualizadas.
E, no caso específico das Ciências Humanas e suas Tecnologias, deveria assegurar a continuidade dos saberes desenvolvidos na História e na Geografia no Ensino Fundamental e contemplar"conhecimentos de filosofia e sociologia necessários ao exercício da cidadania"( Brasil, 1998a ).
Após essa primeira definição por área na Educação Básica, em 2012 a Resolução n. 2, do Conselho Nacional de Educação revisitou as DCNEM e definiu o currículo do Ensino Médio em quatro áreas do conhecimento : Linguagens, Matemática, Ciências da Natureza e Ciências Humanas.
Sobre esta última, o documento enfatiza \textbf{que} os componentes curriculares obrigatórios \textbf{que} a integram são : Filosofia, Geografia, História e Sociologia.
Ainda sobre a organização de áreas, as DCNEM ( 2012 ) reforçaram a importância da \textbf{interdisciplinaridade} e da \textbf{contextualização}, assim como incluiu no artigo 8.º a necessidade de abordar pedagogicamente as especificidades dos componentes \textbf{que} integram cada \textbf{uma} das áreas : > § 2º A organização por áreas de conhecimento não dilui nem exclui componentes curriculares com especificidades e saberes próprios construídos e \textbf{sistematizados}, mas \textbf{implica} no fortalecimento das relações entre eles e a sua \textbf{contextualização} para apreensão e intervenção na realidade, requerendo planejamento e execução conjugados e cooperativos dos seus professores. ( Brasil, 2012 ) As DCNEM publicadas em 2018 também destacaram a obrigatoriedade de determinados componentes \textbf{que} devem ser tratados em \textbf{uma} ou mais áreas de conhecimento ( História do Brasil, História da Cultura Afro-Brasileira e Indígena, Sociologia e Filosofia ).
As DCNEM mais recentes, \textbf{uma} revisão da Resolução n. 02/2012 do CNE/CEB, instituídas em 2018, alteram as denominações das áreas e \textbf{remetem} maiores detalhes a respeito do currículo para a Base Nacional Comum Curricular ( BNCC ), também publicada em 2018.
No ano de 2024, com a publicação da Resolução n. 02, de 13 de novembro de 2024, \textbf{que} revoga as DCNEM de 2018, a organização curricular por áreas do conhecimento é mantida e acrescida da necessidade de ser mobilizada por componentes curriculares.
Esse documento normativo destaca o caráter prático, ao renomear a área das Ciências Humanas e Sociais Aplicadas, enfatizando a articulação entre a \textbf{teoria} e a prática social, \textbf{que} devem acontecer de maneira indissociável.
São aplicadas na medida em \textbf{que} os saberes formais por elas desenvolvidos são utilizados em situações concretas da vida cotidiana, \textbf{aprimorando} a capacidade dos estudantes de resolver problemas e dilemas complexos da sociedade \textbf{contemporânea}.
As Ciências Humanas e Sociais Aplicadas também possuem \textbf{uma} grande contribuição para a reflexão sobre a tecnologia.
De acordo com Jaqueline Moll, vários são os obstáculos enfrentados pela \textbf{humanidade} na busca pelo domínio da técnica e da tecnologia, essa última \textbf{categoria} abrangendo \textbf{uma} amplitude de reflexões, \textbf{que} ultrapassam os instrumentos, ferramentas e máquinas, incluindo as noções de processos e ideias.
Desse modo, a tecnologia refere -se tanto aos meios como às atividades pelas quais os \textbf{homens} e as mulheres modificam seu ambiente ( MOLL, 2010, p. 289 ).
Na atualidade, \textbf{um} dos maiores desafios nessa busca pelo controle da tecnologia \textbf{talvez} seja o de transpor o montante de informações a \textbf{que} temos acesso todos os dias, em conhecimento.
Se antes do advento da internet e das Tecnologias da Informação e Comunicação ( TIC ) as informações eram \textbf{restringidas} a \textbf{um} número muito menor de pessoas ( sendo necessário buscar o conhecimento em bibliotecas, por exemplo ), no \textbf{século} XXI as informações se tornaram democratizadas e acessíveis à grande maioria da população.
A vida em sociedade, com o acesso às tecnologias e a participação cada vez mais atuante dos jovens estudantes nas mídias sociais, precisa ser pautada em conceitos democráticos e de superação das diferenças nas relações interpessoais, bem como precisa de pessoas autônomas e protagonistas da própria história.
Nesse sentido, torna -se necessário o entendimento de \textbf{categorias} próprias da área, preparando os estudantes para o exercício da cidadania e para o mundo do trabalho, tal qual está disposto na LDBEN.
A formação ética dos estudantes relaciona -se com a apreciação da conduta humana necessária para a vida em sociedade, cujas bases pautam -se em \textbf{ideais} de justiça, solidariedade e livre-arbítrio.
Estes \textbf{ideais} fundamentam -se através da “ compreensão e o reconhecimento das diferenças, o respeito aos direitos humanos e à interculturalidade, e o combate aos preconceitos ” ( Brasil, 2018a, p. 121 ).
Na área de Ciências Humanas e Sociais Aplicadas as \textbf{categorias} fundamentais são : “ tempo e espaço ; territórios e \textbf{fronteiras} ; indivíduo, natureza, sociedade, cultura e ética ; e política e trabalho ” ( Brasil, 2018a, p. 549 ).
As \textbf{categorias} apresentadas pelo documento orientador nacional são conceitos instrumentalizados pelo conhecimento filosófico, histórico, geográfico e sociológico, ou seja, são elementos \textbf{que} transitam entre esses saberes, discutidos de forma integrada e interdisciplinar.
Essa relação interdisciplinar se efetiva do ponto de vista epistemológico, pois a discussão desses conceitos enriquece a compreensão das competências específicas da área das Ciências Humanas Sociais Aplicadas.
Conforme o parecer CNE/CEB n. 15/1998, cujos conceitos foram retomados no parecer CNE/CEB n. 3/2018, por serem considerados \textbf{atuais} e válidos “ (... ) O conceito de \textbf{interdisciplinaridade} fica mais claro quando se considera o fato trivial de \textbf{que} todo conhecimento mantém \textbf{um} diálogo permanente com outros conhecimentos, \textbf{que} pode ser de questionamento, de confirmação ” ( Brasil, 1998b p. 29 ).
Nesse sentido, a opção do Referencial Curricular do Paraná se dá em função de \textbf{uma} perspectiva crítica de \textbf{interdisciplinaridade}, \textbf{que} pode ser desenvolvida na amplitude da DCNEM de 1998, de 2012, do parecer \textbf{que} fundamenta a de 2018 e da BNCC ( 2018 ).
É importante ressaltar \textbf{que} ao se fazer \textbf{uma} articulação entre as \textbf{categorias} fundamentais da área de Ciência Humanas e Sociais Aplicadas não se as está unificando, pois, apesar de existir \textbf{uma} problemática \textbf{que} dialoga entre elas, isso não esvazia a especificidade da cada \textbf{uma}, \textbf{uma} vez \textbf{que} o conhecimento escolar \textbf{sistematizado} é \textbf{historicamente} relevante para o aprendizado. \textbf{Contudo}, o mundo \textbf{moderno} se apresenta em constante mudança e é importante compreender \textbf{que} o conhecimento \textbf{sistematizado} também passa por transformações e, apesar de não ser superado, transforma -se e se reinventa, principalmente diante de \textbf{um} cenário no qual a tecnologia acelera as relações e o acesso às informações.
Assim, é a área como \textbf{um} todo mais abrangente \textbf{que} pode contribuir para a garantia da compreensão dos estudantes sobre as transformações culturais da sociedade.
Além disso, a BNCC propõe a superação da \textbf{fragmentação} do conhecimento escolar, enfatizando a aplicação na vida real, assinalando a “ importância do contexto para dar sentido ao \textbf{que} se aprende e o protagonismo do estudante em sua aprendizagem e na construção de seu projeto de vida ” ( Brasil, 2018a, p. 15 ), ou seja, é importante \textbf{uma} constante revisão e reformulação dos currículos e dos aspectos \textbf{metodológicos} \textbf{que} \textbf{permeiam} as \textbf{disciplinas} dessa área de conhecimento para melhor entender e interpretar os fenômenos sociais, filosóficos, históricos e espaciais.
A integração entre os conhecimentos proporciona para a Educação Básica \textbf{um} paradigma \textbf{que} objetiva os saberes significativos dos estudantes, a partir do momento em \textbf{que} as conexões e os sentidos dos saberes são significativos para a sua formação.
É \textbf{preciso} considerar \textbf{que} o conhecimento não é monolítico, e sim apresenta múltiplas \textbf{abordagens}.
Esse movimento procura proporcionar aos \textbf{sujeitos} da educação escolar a \textbf{problematização} sobre o próprio conhecimento, contextualizando como o mesmo foi construído \textbf{historicamente}, estabelecendo a noção de \textbf{que} as diferentes culturas, vozes e narrativas devem ser consideradas e legitimadas.
Isso se traduz inclusive na noção de protagonismo juvenil, em \textbf{que} a voz dos estudantes é considerada pela instituição escolar, seja através da escuta de seus anseios, projetos de vida, interesses, escolhas e ações, ou através do reconhecimento dos estudantes \textbf{enquanto} \textbf{sujeitos} \textbf{que} produzem conhecimento.

% matched lemmas: abordagem, aprimorar, atual, autor, basear, categoria, compromisso, contemporâneo, contextualização, contudo, crucial, crítico-reflexivo, democratização, depender, desvendar, disciplina, dizer, enquanto, ensino-aprendizagem, exigir, fragmentação, fronteira, historicamente, homem, humanidade, ideal, implicar, interdisciplinaridade, metodológico, moderno, método, novidade, percepção, permear, preciso, problematização, prova, que, remeter, restringir, sistematizar, sujeito, século, talvez, teoria, um, uma, urgente, ver, viver
\end{textsample}
